\begin{examplebox}
	\textbf{\textcolor{CExample}{例 1（基础）}}

	\textbf{\textcolor{CExample}{题目：}}
	在 $\triangle ABC$ 中，已知 $a=2$, $A=30^\circ$，求外接圆半径 $R$。

	\textbf{\textcolor{CExample}{解题步骤：}}
	\begin{description}[style=nextline, leftmargin=2.8em, labelsep=0.5em, itemsep=0.45em]
		\item[\textbf{第一步（识别条件）：}] 已知边 $a=2$ 及对角 $A=30^\circ$。
			\par\textbf{理由：} 正弦定理 $R = \frac{a}{2\sin A}$ 需要边长与对角。
		\item[\textbf{第二步（代入公式）：}] 由 $R = \frac{a}{2\sin A}$，得
			\[
				R = \frac{2}{2\sin 30^\circ}.
			\]
			\par\textbf{理由：} 直接套用结论。
		\item[\textbf{第三步（计算）：}] $\sin 30^\circ = \frac12$，故
			\[
				R = \frac{2}{2\times \frac12} = \frac{2}{1} = 2.
			\]
			\par\textbf{理由：} 化简得出数值。
	\end{description}

	\textbf{\textcolor{CExample}{关键结论：}}
	\textbf{$R=2$}。

	\medskip
	\textbf{\textcolor{CExample}{例 2（稍变形）}}

	\textbf{\textcolor{CExample}{题目：}}
	在 $\triangle ABC$ 中，已知 $a=2$, $b=\sqrt{3}$, $A=45^\circ$，求 $R$ 及 $\sin B$。

	\textbf{\textcolor{CExample}{解题步骤：}}
	\begin{description}[style=nextline, leftmargin=2.8em, labelsep=0.5em, itemsep=0.45em]
		\item[\textbf{第一步（求$R$）：}] 由 $R = \frac{a}{2\sin A}$ 得
			\[
				R = \frac{2}{2\sin 45^\circ} = \frac{2}{2\cdot\frac{\sqrt2}{2}} = \frac{2}{\sqrt2} = \sqrt2.
			\]
			\par\textbf{理由：} $\sin 45^\circ = \frac{\sqrt2}{2}$，化简后得 $R = \sqrt2$。

		\item[\textbf{第二步（求$\sin B$）：}] 由正弦定理 $\frac{b}{\sin B} = 2R$ 得
			\[
				\sin B = \frac{b}{2R} = \frac{\sqrt3}{2\sqrt2} = \frac{\sqrt6}{4}.
			\]
			\par\textbf{理由：} $b$ 与对角 $B$ 满足同一比值 $2R$，代入 $R$ 即可。

		\item[\textbf{第三步（验证）：}] 检查 $\sin B = \frac{\sqrt6}{4} \approx 0.612 < 1$，三角形唯一（$B$ 为锐角）。
			\par\textbf{理由：} 保证三角形存在且解有效。
	\end{description}

	\textbf{\textcolor{CExample}{关键结论：}}
	\textbf{$R = \sqrt2$，$\sin B = \frac{\sqrt6}{4}$}。

	\medskip
	\textbf{\textcolor{CExample}{例 3（综合应用）}}

	\textbf{\textcolor{CExample}{题目：}}
	在 $\triangle ABC$ 中，外接圆半径 $R=1$，且 $\angle A = 60^\circ$，$\angle B = 45^\circ$，求 $\triangle ABC$ 的周长。

	\textbf{\textcolor{CExample}{解题步骤：}}
	\begin{description}[style=nextline, leftmargin=2.8em, labelsep=0.5em, itemsep=0.45em]
		\item[\textbf{第一步（求$\sin$值）：}] 计算 $\sin 60^\circ = \frac{\sqrt3}{2}$，$\sin 45^\circ = \frac{\sqrt2}{2}$。
			\[
				\sin C = \sin(180^\circ - 60^\circ - 45^\circ) = \sin 75^\circ = \frac{\sqrt6 + \sqrt2}{4}.
			\]
			\par\textbf{理由：} 三角形内角和 $180^\circ$，已知两角可求第三角；$\sin 75^\circ$ 需用和角公式或熟记。

		\item[\textbf{第二步（由$R$求各边）：}] 由正弦定理 $a = 2R\sin A$，$b = 2R\sin B$，$c = 2R\sin C$，代入 $R=1$ 得
			\[
				a = 2\sin 60^\circ = \sqrt3,
				\qquad
				b = 2\sin 45^\circ = \sqrt2,
			\]
			\[
				c = 2\sin 75^\circ = \frac{\sqrt6 + \sqrt2}{2}.
			\]
			\par\textbf{理由：} $2R\sin \theta$ 是正弦定理的变形。

		\item[\textbf{第三步（求周长）：}] 周长 $L = a+b+c$，代入得
			\[
				L = \sqrt3 + \sqrt2 + \frac{\sqrt6 + \sqrt2}{2} = \sqrt3 + \frac{3\sqrt2}{2} + \frac{\sqrt6}{2}.
			\]
			\par\textbf{理由：} 合并同类项。
	\end{description}

	\textbf{\textcolor{CExample}{关键结论：}}
	\textbf{周长 $L = \sqrt3 + \frac{3\sqrt2}{2} + \frac{\sqrt6}{2}$}。
\end{examplebox}
